\documentclass[10pt,a4paper]{amsart}
\usepackage[margin=19mm]{geometry}
\usepackage{amsmath,amssymb,mathtools}
\usepackage[scr=rsfs,cal=euler]{mathalfa}
\usepackage{xfrac}
\usepackage[unicode,hidelinks,bookmarks=false]{hyperref}
\newcommand{\ie}{i.e.\ }
\newcommand{\cf}{cf.\ }
\newcommand{\resp}{respectively\ }
\newcommand{\Subsection}{\subsection}
\providecommand{\nobreakdash}{\nobreak\hbox{-}}
\setlength{\emergencystretch}{2em}
\multlinegap0pt
\usepackage{bm}
\usepackage{bbm}
\usepackage{dsfont}
\usepackage{xspace}
\usepackage{cancel}
\usepackage{mathtools}
\usepackage{amsthm}
\makeatletter
\@ifundefined{lemma}{\newtheorem{lemma}{Lemma}}{}
\makeatother
\def \OC#1{\text{\bf O}(#1,\cn)}
\def \SOC#1{\text{\bf SO}(#1,\cn)}
\def \cn{{\mathbb C}}
\def \rn{{\mathbb R}}
\title[Complete Minimal Submanifolds of the Non-Compact Duals to th]{Complete Minimal Submanifolds of the Non-Compact Duals to the Classical Compact Lie Groups}
\author{Thomas Jack Munn}
\date{Source: arXiv:2608.03354v1 (CC BY 4.0)}
\begin{document}
\maketitle

\noindent\emph{Source and licence.} Complete Minimal Submanifolds of the Non-Compact Duals to the Classical Compact Lie Groups. Version arXiv:2608.03354v1 (CC BY 4.0), distributed under the Creative Commons Attribution 4.0 International licence, as recorded in the arXiv licence metadata for this exact version and shown on the abstract page https://arxiv.org/abs/2608.03354v1. Full rights notice: licenses/arxiv-2608.03354-CC-BY-4.0.txt.\par\smallskip

\noindent\textit{Editorial scope.} Counted unit: the Lemma whose source label is \texttt{lemma-zeros-sonC} in the article (source0.tex lines 429--430, source label \texttt{lemma-zeros-sonC}), delivered together with the article's own complete original proof of it (source0.tex lines 431--486, one \texttt{proof} environment of 56 lines). No mathematics is written, repaired or re-derived: statement and proof are the article's own text line for line. Required context retained uncounted: none. Every definition, display and label of those lines is retained unchanged. Omitted from this package and not redistributed: 1--428, 487--1001. Nothing inside a retained excerpt is omitted or altered: every retained body is delivered exactly as the article's lines are written, so no omission receipt of any kind accompanies this package. Citations: the retained excerpts carry no citation command at all, so no bibliography entry of the article is transcribed and no reference is redistributed. Reference adaptation: none was needed and none was made; every reference inside the delivered unit resolves inside it, so no unresolved reference or citation is produced. Printed slips kept literal: no printed word, symbol, hypothesis or number of the counted statement or of its proof was altered. Layout and class: the article's own class is \texttt{amsart} with packages including color, amsmath, amsfonts, amscd, amssymb, url, hyperref, babel, booktabs, tikz; the wrapper is the lane's pinned \texttt{amsart} editorial wrapper at 10pt on A4, which restates the theorem-like environments the retained text needs and adds the article's own macro definitions that the retained text needs (\texttt{\textbackslash OC}, \texttt{\textbackslash SOC}, \texttt{\textbackslash cn}, \texttt{\textbackslash rn}), with extra wrapper packages (bm, bbm, dsfont, xspace, cancel, mathtools). The delivered numbering is the unit's own. Assets: no retained span contains a figure environment, a graphics-inclusion command, a PDF-page inclusion, a table or a TikZ picture; no image file is copied into this package, the package declares no asset dependency and no image file is redistributed with it. Licence: version arXiv:2608.03354v1 of this article is distributed under the Creative Commons Attribution 4.0 International licence, as recorded in the arXiv licence metadata for this exact version and shown on the abstract page https://arxiv.org/abs/2608.03354v1. A full rights notice is delivered at licenses/arxiv-2608.03354-CC-BY-4.0.txt. Characters: every retained excerpt is pure ASCII, so the delivered engine, XeLaTeX with its default Latin Modern text font, reports no missing character.
\begin{lemma}\label{lemma-zeros-sonC}
	Let $u,v \in \mathbb{C}^n$ be linearly independent vectors for $n \geq 3$. Then there exists $g \in \SOC{n}$ such that $gu$ and $gv$ are orthogonal with respect to the standard Hermitian inner product.\end{lemma} 

\begin{proof}
	First recall that $\OC{n}$ is the isometry group of the complex bilinear inner product on $\cn^n$, so by Witt's theorem it is sufficient to define an isometry between $\text{span}\{u,v\}$ and its image $\text{span}\{x=gu,y=gv\}$ with the property that $x$ and $y$ are Hermitian-orthogonal.	
	
	 For notational convenience we only display the first coordinates of the vectors in $\cn^n$, as the other entries are all zero.
	 
	Now we show that such $x,y$ always exist. Let 
	\begin{eqnarray*}
		&a = u^t u, \ b = v^t v \text{ and } c = u^t v.&
	\end{eqnarray*}
	
	We will construct $x,y$ satisfying
	\begin{eqnarray}\label{eqn-xy}
		&a = x^t x, \ b = y^t y, \ c = x^t y& \text{and } x^*y=0.
	\end{eqnarray}
	We proceed by cases. \\

	
	First suppose that $a = b =c = 0$. Note that this case can only occur when $n \geq 4$, since the maximal isotropic subspace in $\cn^3$ is one-dimensional. Let

$$ x = \begin{pmatrix}
		1 \\ i \\ 0 \\ 0
	\end{pmatrix} \text{ and }  \ y = \begin{pmatrix}
		0 \\ 0 \\ 1 \\ i
	\end{pmatrix}$$
	which clearly satisfy equations (\ref{eqn-xy}).
	
	Now consider the case when $a = 0$ but $(b,c) \neq (0,0)$. Then choose $s \in \cn$ with $s^2=b$ and define
	$$ x = \begin{pmatrix}
		1 \\ i \\ 0
	\end{pmatrix} \text{ and }  \ y = \begin{pmatrix}
		\tfrac{c}{2} \\ -\tfrac{ic}{2}\\ s
	\end{pmatrix}.$$
A simple calculation shows $x$ and $y$ satisfy equations (\ref{eqn-xy}) and since $s \neq 0$ or $c \neq 0$, $x$ and $y$ are clearly linearly independent. By symmetry, an identical argument works when $b=0$ and $(a,c) \neq (0,0)$.

	Finally we need to consider the case when neither $a$ nor $b$ are zero. Choose $r \in \cn$ with $r^2 = a$ and $t \in \rn \setminus \{0\}$ then let 
	
		$$ x = r \cdot \begin{pmatrix}
		\cosh t \\ i \sinh t \\ 0
	\end{pmatrix} \text{ and }  \ \tilde y = \begin{pmatrix}
		i\sinh t\\ \cosh t\\ 0
	\end{pmatrix}.$$ 
	A computation shows
 \begin{eqnarray*}
 	& x^t x = r^2 = a, \  \tilde y^t \tilde y =1, \  x^t \tilde y= 2i r \cosh t \sinh t = i r \sinh 2t&, \ x^* \tilde y=0.
 \end{eqnarray*}
 Since neither $r$ nor $t$ are zero we have that $x^t \tilde y \neq 0$ so we can pick some $s \in \cn$ such that $s^2 = b-\frac{c^2}{(x^t \tilde y)^2}$ and finally define
 $$ y = \tfrac{c}{x^t \tilde y} \cdot \tilde y + s \cdot \begin{pmatrix}
 	0 \\ 0 \\ 1
 \end{pmatrix}.$$
 It is then clear that $x$ and $y$ satisfy equations (\ref{eqn-xy}).
 
 So, by Witt's theorem there always exists a $G \in \OC{n}$ with the required properties. If $\det G =1$ let $g = G$. Otherwise, if $\det G = -1$ let $R = \text{diag}(-1,1,\dots ,1) \in \OC{n}$ and set $g = RG \in \SOC{n}$. Since $R$ is unitary we have that
	$$ (RGu)^*(RGv) = (Gu)^*R^*R(Gv) =(Gu)^*Gv=0.$$

 
\end{proof}

\end{document}
